\chapter{Logger 与证据台账规范}
\label{chap:workflow}

本章把代码实现路径进一步约束为可投稿证据路径。NeurIPS 主会论文不能只依赖“某个脚本曾经跑出过表格”；它需要每个结果都能回答四个问题：使用了哪一次运行，运行时仓库处于什么状态，输入数据和 split 的 checksum 是什么，统计检验如何从原始 per-query 结果生成。Logger 的职责就是把这些问题变成强制字段。

\begin{table}[H]
	\centering
	\caption{NeurIPS run logger 的最小字段规范。}
	\label{tab:logger-schema}
	\scriptsize
	\resizebox{\textwidth}{!}{%
	\begin{tabular}{p{0.20\textwidth}p{0.16\textwidth}p{0.25\textwidth}p{0.31\textwidth}}
		\toprule
		字段 & 类型 & 角色 & 说明 \\
		\midrule
		\texttt{run\_id} & string & 主键 & 采用日期、实验编号、方法和短 hash 组成，所有 CSV/JSON/MD 共享。 \\
		\texttt{supervisor\_gate} & enum & 监督阶段 & 取值如 E0-freeze、E1-baseline、E2-ablation、E3-leakage。 \\
		\texttt{git\_commit} / \texttt{git\_dirty} & string/bool & 代码状态 & 记录 HEAD、branch、未提交变更摘要；dirty run 只能作探索，不能直接进主表。 \\
		\texttt{command} / \texttt{env} & string/object & 复现入口 & 保存完整命令、Python 版本、关键依赖版本、设备信息。 \\
		\texttt{dataset\_checksum} & string & 数据快照 & 指向 dataset manifest、split manifest 和原始 JSON/CSV checksum。 \\
		\texttt{split\_id} & string & 协议边界 & 明确 204 audio-grouped、legacy 211 或更严格 family split。 \\
		\texttt{query\_count} / \texttt{kb\_count} & integer & 规模审计 & 防止不同口径的样本规模被混入同一主表。 \\
		\texttt{method\_config} & object & 方法定义 & 保存 retriever、embedding、normalization、top-k、seed、ablation flags。 \\
		\texttt{metrics} / \texttt{stats} & object & 统计结果 & 保存均值、CI、paired permutation、Holm correction 和有效样本交集。 \\
		\texttt{artifacts} & list & 证据索引 & 列出 per-query CSV、summary JSON、stats MD、stderr/stdout 和 checksum。 \\
		\texttt{failure\_flags} & list & 负面证据 & 记录缺失缓存、跳过样本、异常输入、退化场景失败类型。 \\
		\bottomrule
	\end{tabular}}
\end{table}

表~\ref{tab:logger-schema} 是后续实验代码需要满足的最小审计接口。它的主要作用是让论文表格、附录统计和匿名代码包共享同一证据主键。

建议的运行产物目录为 \texttt{Experiments/NeurIPSRuns/<run\_id>/}，其中至少包含 \texttt{run.json}、\texttt{per\_query.csv}、\texttt{summary.json}、\texttt{stats.md}、\texttt{checksums.sha256} 和 \texttt{stdout.log}。本报告不声称该目录已经完全实现；它是从当前 evidence boundary 推导出的 design intent。后续实现时，旧的 Protocol-A/Protocol-C 产物可以被迁移或注册到同一 schema，而不需要重写历史结果。

KaiMing 对 logger 的验收标准很简单：任何进入主文的表格行，都必须能反查到唯一 \texttt{run\_id}；任何显著性结论，都必须能反查到配对样本交集；任何缺失输出或 dirty run，都必须在表注或局限中被披露。若做不到这一点，该结果只能进入探索记录或附录待复现清单。

代码库中存在产品链路与实验链路两套路径。产品链路从用户文本和可选音频出发，通过 DeepSeek 解析风格与参数，经 \texttt{import\_params.json} 传入 JUCE 插件，再由插件加载参数和音频并离线渲染。该链路说明 Audio-Agent 具备可执行插件控制的工程背景，但本文不把它作为主实验结果来源，因为主实验没有评价完整插件输出音频质量。

实验链路更适合作为论文证据。数据加载由 \texttt{Experiments/common/dataset\_loader.py} 负责，优先读取外部 1267 条数据集 JSON，再回退到 repo 内 JSON 或 SQLite 合并。TRR 检索由 \texttt{Experiments/common/trr\_adapter.py} 实现，优先使用缓存向量。直接检索比较由 \texttt{Experiments/AblationStudies/direct\_retrieval\_comparison.py} 驱动，Protocol-C 由 \texttt{Experiments/AblationStudies/robustness\_test.py} 驱动，评价指标由 \texttt{Experiments/common/evaluate.py} 计算。

\begin{figure}[H]
	\centering
	\resizebox{\textwidth}{!}{%
	\begin{tikzpicture}[node distance=1.25cm and 1.0cm]
		\node[reportnode, minimum width=3.2cm] (dataset) {数据集 JSON\\1267 records};
		\node[reportprocess, right=of dataset, minimum width=3.0cm] (split) {切分审计\\204/1063};
		\node[reportprocess, right=of split, minimum width=3.2cm] (retrievers) {检索器\\TRR/Text/Wav2Vec/CLAP};
		\node[reportprocess, right=of retrievers, minimum width=2.8cm] (metrics) {参数评价\\Evaluator};
		\node[reportnode, right=of metrics, minimum width=3.2cm] (artifacts) {CSV/JSON/MD\\证据产物};
		\draw[reportarrow] (dataset) -- (split);
		\draw[reportarrow] (split) -- (retrievers);
		\draw[reportarrow] (retrievers) -- (metrics);
		\draw[reportarrow] (metrics) -- (artifacts);
		\node[reportnode, below=1.0cm of retrievers, minimum width=4.2cm, fill=ReportSoftYellow] (boundary) {证据治理\\主证据 / 诊断 / 待补 / 历史};
		\draw[reportarrow] (artifacts) |- (boundary);
	\end{tikzpicture}}
	\caption{实验证据生成路径。论文主结论应回到可复核的 CSV、JSON、Markdown 统计产物，避免依赖未验证的工程愿景。}
	\label{fig:experiment-evidence-flow}
\end{figure}

图~\ref{fig:experiment-evidence-flow} 对应当前最可靠的复现路径。对于 Protocol-A，复核者应从 split audit、audio-grouped per-query CSV 和 objective stats 开始，确认样本规模、指标定义和配对统计。对于 Protocol-C，复核者应检查每个 scenario 下 Text-only、TRR-only 和 Fusion 的 L2、权重均值及 Holm-corrected comparisons。对于 P0，复核者只能先阅读 runbook 和脚本，因为报告引用的原始 outputs 尚未提交。

当前代码中还存在一些不应被误写为主证据的部分。实验侧 Text-RAG 实际更接近 deterministic keyword overlap；Protocol-B 已经被标注为 deprecated；E0 合成音频管线仍有缺失音频记录。因此，本文在结论中使用 “under the current held-out protocol”、“cached TRR representation” 和 “currently evaluated baselines” 等限定语。
